\section{An obstruction to deterministic selection}\label{sec:splitting}

We isolate the measure-theoretic step. A homeomorphism onto an open
image sends Borel sets to Borel sets, since its inverse is
continuous.

\begin{lemma}\label{lem:splitting}
Let \(d\ge2\) be an integer, let \(B\subset\R^d\) be open, let
\(\nu\) be a Borel probability measure with \(\nu(B)=1\), and let
\(H_1,\ldots,H_4:B\to\R^d\) be homeomorphisms onto open images.
Suppose that \(B,H_1(B),\ldots,H_4(B)\) are pairwise disjoint.
Define
\begin{equation}\label{eq:split-measure}
 \mu=\frac13\nu+\frac16\sum_{j=1}^4(H_j)_\#\nu.
\end{equation}
There is no Borel map \(T:\R^d\to\R^d\) such that
\(T_\#\mu=\mu\) and
\begin{equation}\label{eq:selection}
 T(x)\in\{H_1(x),\ldots,H_4(x)\}
 \quad\text{for \(\nu\)-almost every }x\in B.
\end{equation}
\end{lemma}

\begin{proof}
For each \(j\), define the Borel set
\[
 S_j=\{x\in B:T(x)=H_j(x)\}.
\]
Its image \(H_j(S_j)\) is Borel. Pairwise disjointness of the
component images and injectivity of \(H_j\) give
\[
 \mu(S_j)=\frac13\nu(S_j),\qquad
 \mu(H_j(S_j))=\frac16\nu(S_j).
\]
Since \(S_j\subset T^{-1}(H_j(S_j))\), preservation of \(\mu\)
would imply
\[
 \frac13\nu(S_j)
 \le \mu(T^{-1}(H_j(S_j)))
 =\mu(H_j(S_j))
 =\frac16\nu(S_j).
\]
Thus \(\nu(S_j)=0\) for all four indices, contradicting
\eqref{eq:selection}.
\end{proof}

The argument does not require \(\nu\) to be uniform or even
absolutely continuous. Additional points mapped into \(H_j(S_j)\)
can only increase its preimage mass. Matching the total mass of
each component by a partition therefore cannot remove the
obstruction.

